\documentclass[11pt]{article}
\usepackage{amsmath,amssymb,amsthm}
\usepackage[margin=1in]{geometry}
\usepackage{booktabs}
\usepackage{hyperref}
\usepackage{xcolor}
\hypersetup{colorlinks=true, linkcolor=blue!50!black, urlcolor=blue!50!black}

\newcommand{\Var}{\mathcal{V}}
\newcommand{\code}[1]{\texttt{#1}}
\DeclareMathOperator{\Sol}{Sol}
\DeclareMathOperator{\GCRD}{GCRD}
\DeclareMathOperator{\ord}{ord}
\DeclareMathOperator{\lclm}{LCLM}

\newtheorem{definition}{Definition}
\newtheorem{lemma}{Lemma}
\newtheorem{proposition}{Proposition}
\newtheorem{theorem}{Theorem}
\newtheorem{example}{Example}
\newtheorem{remark}{Remark}

\title{Recovering the partial-solution strata by membership alone:\\
       a GCRD closure construction}
\author{}
\date{\today}

\begin{document}
\maketitle

\noindent
A companion note to \emph{A New Method of Solving PDEs}, and a sequel to the
note \emph{Membership versus variety, and the partial-solution strata}. That
note established that the core (membership) algorithm computes the locus
$\Var = \{c : P \in [A(c)]\}$ --- a \emph{universal} condition --- and is
therefore blind, by construction, to the \emph{partial-solution strata}:
parameter values $c$ at which a proper nontrivial subfamily of the ansatz
solutions satisfies the PDE while the family as a whole does not. Capturing
those strata seemed to require adjoining $P$ to the ansatz and running a full
differential decomposition --- the expensive existential computation.

This note shows that, for ans\"atze whose top element is a \emph{linear} ODE,
the existential computation is unnecessary. Given such an ansatz $A$, we
construct --- by pure polynomial linear algebra, with no differential
elimination beyond the Ritt reduction the core algorithm already performs ---
a finite set of auxiliary ans\"atze $A_1, \dots, A_{n-1}$ of the same class
and lower order, such that every partial-solution stratum of $A$ is an
ordinary \emph{membership} stratum of some $A_k$. Running the core algorithm
on the collection then achieves the completeness of the existential
formulation at membership-formulation cost. The engine is the greatest common
right divisor (GCRD) of Ore polynomials, made parametric through the
subresultant theory of Chardin and Li \cite{chardin,li}; as a by-product, the
partial strata themselves acquire explicit polynomial defining equations in
the \emph{original} constant space.

\section{The covering specification}

Throughout, $A(c)$ is an ansatz with constant parameters
$c \in \mathbb{C}^{n_c}$, $P$ is the PDE, and
$\Var(A) = \{c : P \in [A(c)]\}$ is the membership locus computed by the core
algorithm (Ritt reduction of $P$ modulo $A$, projection of the remainder onto
the constants).

\begin{definition}[$\forall$-cover]
\label{def:cover}
A finite collection of ans\"atze $\{A_1, \dots, A_m\}$ (each with its own
constants) \emph{$\forall$-covers} $A$ with respect to $P$ if, for every
$c^*$ and every nonzero
$\Psi_0 \in \Sol(A(c^*)) \cap \Sol(P)$, there exist an index $i$ and
constants $c'$ with
\[
\Psi_0 \in \Sol\big(A_i(c')\big)
\qquad\text{and}\qquad
c' \in \Var(A_i).
\]
\end{definition}

If a collection $\forall$-covers $A$, then running the core algorithm on each
member delivers every solution function that the existential computation on
$A$ would deliver --- each one packaged inside a family that solves $P$
wholesale, which is the form of answer the method is designed to produce. Two
requirements make the notion useful rather than vacuous:

\begin{enumerate}
\item The collection must be constructed from $A$ and the \emph{syntactic
      shape} of $P$ (its order and degree bounds) alone. If the construction
      required decomposing $A \cup \{P\}$, it would cost what it saves.
\item The collection must be finite, with all degree bounds computable.
\end{enumerate}

\section{Setting: a single linear ODE element}
\label{sec:setting}

We treat the shape of the paper's Figure ``ansatz 5'': a base differential
ring $B$ (e.g.\ $\mathbb{Q}\{x,y,z\}[r]/(r^2 - x^2 - y^2 - z^2)$), a linear
form $v = v_1 x + v_2 y + v_3 z + v_4 r$ with constant coefficients, the
chain-rule relations $\Psi_x = \Psi' v_x$, etc., and a top element
\begin{equation}
\label{eq:L}
L \;=\; \sum_{i=0}^{n} a_i(v)\,\partial^i,
\qquad a_i \in K[v],\ \deg_v a_i \le d,\ a_n \neq 0,
\end{equation}
regarded as an Ore (differential) operator acting on $\Psi$ as a function of
$v$; here $K = \mathbb{Q}(c)$ and $\partial = d/dv$. The natural home for the
operator computations is the right-Euclidean ring $K(v)[\partial]$, in which
every left ideal is principal and greatest common right divisors exist. Two
standard facts \cite{ore,vdps} are used repeatedly:

\begin{enumerate}
\item ($G$ exists and is rational.) For operators $L, R_1, \dots, R_s \in
      K(v)[\partial]$, the left ideal $K(v)[\partial]L + \sum_\beta
      K(v)[\partial]R_\beta$ is principal; its monic generator
      $G = \GCRD(L, R_1, \dots, R_s)$ right-divides each argument and is
      computable by the right Euclidean algorithm, with coefficients in
      $K(v)$.
\item (Kernels intersect.) On a domain where the coefficients are analytic
      and $a_n \neq 0$, the local solution space of an order-$m$ operator is
      $m$-dimensional, and
      \[
      \ker G \;=\; \ker L \,\cap\, \bigcap_\beta \ker R_\beta .
      \]
\end{enumerate}

\paragraph{The remainder as a system of operators.}
Let $P$ be a differential polynomial \emph{linear} in $\Psi$ and its
derivatives. Ritt's full reduction of $P$ modulo $A$ yields
$h P - r \in [A]$ with $r$ fully reduced; since both $P$ and the elements of
$A$ are linear in the $\Psi$-jets, so is $r$, and full reducedness confines
it to $\Psi, \Psi', \dots, \Psi^{(n-1)}$. Its coefficients lie in the base
ring extended by $v$. Expanding along a set of power products
$m_\beta$ in the base indeterminates that remain linearly independent over
$K(v)$ on the solution locus (Remark~\ref{rem:splitting}), we write
\begin{equation}
\label{eq:split}
r \;=\; \sum_\beta m_\beta \cdot R_\beta(\Psi),
\qquad R_\beta \in K(v)[\partial],\quad \ord R_\beta \le n - 1 .
\end{equation}
In the univariate case (base ring $\mathbb{Q}\{x\}$, $v = x$) there is a
single $\beta$ and $R_1 = r$ itself.

\begin{remark}[base-splitting hypothesis]
\label{rem:splitting}
Because $v$ is itself a linear form in the base variables, monomials in
$x, y, z, r$ are \emph{not} automatically independent over $K(v)$: the
relation $v_1 x + v_2 y + v_3 z + v_4 r - v = 0$ is a dependence. A valid
splitting eliminates one base coordinate in favor of $v$ (possible whenever
the form $v$ is nondegenerate, a genericity condition on the $v_i$ of the
same character as the paper's hypothesis~(2)) and expands $r$ in the standard
monomials of the remaining coordinates. We assume such a splitting has been
fixed; all statements below are relative to it.
\end{remark}

\begin{lemma}[solution criterion]
\label{lem:criterion}
Let $c^*$ satisfy $h(c^*) \neq 0$. For every solution $\Psi_0$ of $A(c^*)$:
\[
P(\Psi_0) = 0
\quad\Longleftrightarrow\quad
R_\beta(c^*)\,\Psi_0 = 0 \ \text{ for all } \beta .
\]
\end{lemma}

\begin{proof}
($\Leftarrow$) All $R_\beta(c^*)\Psi_0 = 0$ gives $r(c^*)(\Psi_0) = 0$
by \eqref{eq:split}; the reduction identity gives
$h(c^*)\,P(\Psi_0) = r(c^*)(\Psi_0) = 0$, and $h(c^*) \neq 0$.
($\Rightarrow$) $P(\Psi_0) = 0$ forces $r(c^*)(\Psi_0) \equiv 0$ on the
domain; by the independence of the $m_\beta$ over functions of $v$
(Remark~\ref{rem:splitting}), each coefficient function
$R_\beta(c^*)\Psi_0$ vanishes.
\end{proof}

\section{The GCRD structure theorem}

Fix $c^*$ with $h(c^*) \neq 0$ and set
\[
G(c^*) \;=\; \GCRD\big(L(c^*),\, R_1(c^*), \dots, R_s(c^*)\big)
\;\in\; \mathbb{C}(v)[\partial],
\qquad k \;=\; \ord G(c^*) .
\]
Since $G$ right-divides $L$, we have $0 \le k \le n$. Write
$\widetilde{G}$ for $G$ with denominators cleared (polynomial coefficients,
content removed), and let $A_G(c^*)$ denote the ansatz $A(c^*)$ with the top
element $L(c^*)$ replaced by $\widetilde{G}(c^*)$ --- the same tower, one
order lower.

\begin{theorem}[trichotomy and refinement]
\label{thm:structure}
With the hypotheses of Lemma~\ref{lem:criterion}:
\begin{enumerate}
\item $k = n$ iff every $R_\beta(c^*) = 0$, iff $c^* \in \Var(A)$
      (membership: the whole family solves $P$).
\item $k = 0$ iff $\Sol(A(c^*)) \cap \Sol(P) = \{0\}$ (no stratum).
\item $0 < k < n$ iff $c^*$ lies on a partial-solution stratum, and then
      \[
      \Sol(A(c^*)) \cap \Sol(P) \;=\; \Sol\big(A_G(c^*)\big),
      \]
      a $k$-dimensional family: exactly the solutions of the order-$k$
      operator $G(c^*)$, transported through the unchanged lower tower.
\end{enumerate}
\end{theorem}

\begin{proof}
By Lemma~\ref{lem:criterion} and fact (2) of \S\ref{sec:setting},
$\Sol(A(c^*)) \cap \Sol(P)$ consists of the solutions of the tower whose top
function lies in $\ker L \cap \bigcap_\beta \ker R_\beta = \ker G$. If every
$R_\beta(c^*) = 0$ the left ideal is generated by $L$ alone, $G \sim L$ and
$k = n$; the projection coefficients of $r$ vanish at $c^*$, which is the
definition of $\Var(A)$. If $k = 0$ then $G$ is a unit and the common kernel
is trivial. Otherwise $\ker G$ is $k$-dimensional with $0 < k < n$: a proper
nontrivial subfamily, i.e.\ a partial stratum in the sense of the previous
note's Definition~1; and since $G$ right-divides $L(c^*)$, every solution of
$A_G(c^*)$ solves $A(c^*)$, giving the displayed equality.
\end{proof}

The next statement is the reason the construction pays: the refined system is
not merely consistent with $P$ --- it \emph{members} $P$, so the core
algorithm, run on a generic family containing $A_G$, finds it.

\begin{theorem}[membership transfers to the refinement]
\label{thm:membership}
Let $0 < k < n$ as in Theorem~\ref{thm:structure}(3), and let $A_k$ be the
\emph{generic} ansatz obtained from $A$ by replacing the top element with an
undetermined linear ODE of order $k$,
\[
L_k \;=\; \sum_{j=0}^{k} b_j(v)\,\partial^j,
\qquad \deg_v b_j \le B_k,
\]
whose coefficients' $v$-coefficients form the new constant vector $c'$
(the bound $B_k$ is fixed in Proposition~\ref{prop:bound}). Let
$c'(c^*)$ be the coefficient vector of $\widetilde{G}(c^*)$. Then
\[
c'(c^*) \;\in\; \Var(A_k),
\]
i.e.\ the Ritt remainder of $P$ modulo $A_k$, projected onto the constants,
vanishes at $c'(c^*)$.
\end{theorem}

\begin{proof}
Reduce $P$ modulo $A_k$ over $K(c')$: $h_k P - r_k \in [A_k]$ with $r_k$
linear in $\Psi, \dots, \Psi^{(k-1)}$, and split
$r_k = \sum_\beta m_\beta R_{k,\beta}$ as in \eqref{eq:split}, so each
$R_{k,\beta} \in K(c')(v)[\partial]$ has order $\le k - 1$. Specialize
$c' = c'(c^*)$; note $h_k(c'(c^*)) \neq 0$, since the initial and separant of
the monic-in-$\Psi^{(k)}$ element $\widetilde{G}$ reduce to its leading
$v$-coefficient, a nonzero polynomial. Every
$\Psi_0 \in \Sol(A_k(c'(c^*))) = \Sol(A_G(c^*))$ solves $P$ by
Theorem~\ref{thm:structure}(3), so Lemma~\ref{lem:criterion}, applied to
$A_k$, gives $R_{k,\beta}(c'(c^*))\,\Psi_0 = 0$ for all $\beta$ and all
$\Psi_0$ in a $k$-dimensional space. An operator of order $\le k - 1$
annihilating $k$ linearly independent functions (nonzero Wronskian) is the
zero operator. Hence every $R_{k,\beta}(c'(c^*)) = 0$; equivalently, every
projection coefficient of $r_k$ vanishes at $c'(c^*)$, which is
$c'(c^*) \in \Var(A_k)$.
\end{proof}

Note the Wronskian argument replaces the characteristic-set machinery of the
paper's Theorem~3 entirely --- linearity is doing all the work.

\section{Parametric stratification and the degree bound}
\label{sec:subresultants}

Theorems~\ref{thm:structure} and~\ref{thm:membership} are pointwise in
$c^*$. The subresultant theory for Ore polynomials \cite{chardin,li} makes
them parametric. For a pair $L, R$ of orders $n > n'$, the $k$-th
subresultant $S_k(L, R)$, $0 \le k \le n'$, is an operator of order $k$
whose coefficients are determinants of Sylvester-type matrices of size
$n + n' - 2k$ built from the coefficients of
$\partial^j R$ ($0 \le j < n - k$) and $\partial^j L$ ($0 \le j < n' - k$);
as in the commutative theory,
\[
\ord \GCRD(L, R) \;=\; \min\{\,k : S_k \not\equiv 0\,\},
\qquad
\GCRD(L, R) \;\sim\; S_{k_{\min}} .
\]
Since Ritt's full reduction is fraction-free, the coefficients of $L$ and of
every $R_\beta$ are polynomials in $c$ and $v$; hence so is every entry of
every subresultant matrix, and we obtain:

\begin{proposition}[strata equations]
\label{prop:strata}
For each $k$, the locus $\{c : \ord G(c) \ge k\}$ is closed, cut out by the
vanishing of the $v$-coefficients of the subresultants
$S_0, \dots, S_{k-1}$ --- explicitly computable polynomials in $c$. In
particular the union of the partial-solution strata is the constructible set
$\{\ord G \ge 1\} \setminus \Var(A)$, now presented by polynomial equations
and inequations in the \emph{original} constants, refining the inclusion
$\pi_c(\Sol(A) \cap \Sol(P)) \supseteq \Var$ of the previous note into a
computable stratification. Moreover, on the stratum $\{\ord G = k\}$ the
coefficients of $\widetilde{G}$ --- hence the transfer map
$c \mapsto c'(c)$ of Theorem~\ref{thm:membership} --- are polynomial in $c$.
\end{proposition}

\begin{proposition}[degree bound]
\label{prop:bound}
Let $D$ bound the $v$-degrees of the coefficients of $L$ and of all
$R_\beta$, and $n' = \max_\beta \ord R_\beta \le n-1$. Then the
$v$-degrees of the coefficients of $S_k$ --- and hence of
$\widetilde{G}$ on any stratum where $\ord G = k$ --- are at most
\[
B_k \;=\; (n + n' - 2k)\, D .
\]
\end{proposition}

\begin{proof}[Proof sketch]
Left-multiplying by $\partial$ maps an operator with polynomial coefficients
of $v$-degree $\le D$ to one with coefficient degrees $\le D$
(differentiation does not raise polynomial degree). Every entry of the
$k$-th matrix is a coefficient of some $\partial^j L$ or $\partial^j R_\beta$
and so has degree $\le D$; the determinant of a square matrix of size
$n + n' - 2k$ then has degree $\le (n + n' - 2k) D$.
\end{proof}

The bound is coarse but uniform in $c$, which is what the construction
needs. (For $s > 1$ operators $R_\beta$, compute the GCRD by pairwise
iteration, $\GCRD(L, R_1, R_2) = \GCRD(\GCRD(L, R_1), R_2)$, re-clearing
denominators at each step; the bounds compose and remain computable, at the
cost of an explicit closed formula.)

\section{The algorithm}
\label{sec:algorithm}

\noindent\textbf{Input:} an ansatz $A$ of the shape of \S\ref{sec:setting}
(top element linear of order $n$, coefficient degrees $\le d$); the order and
degree shape of the PDE $P$.\\
\textbf{Output:} ans\"atze $A_1, \dots, A_{n-1}$ such that
$\{A, A_1, \dots, A_{n-1}\}$ $\forall$-covers $A$ off the bad locus;
optionally, the partial strata as subvarieties of the original $c$-space.

\begin{enumerate}
\item \textbf{Reduce.} Ritt-reduce $P$ modulo $A$ over $\mathbb{Q}(c)$
      (this is the core algorithm's own first step), obtaining $h$ and $r$;
      split $r$ over base power products as in
      equation~\eqref{eq:split} to get the operators $R_\beta$.
\item \textbf{Stratify.} For $k = 1, \dots, n-1$ compute the subresultant
      data of $L$ against $\{R_\beta\}$: the strata polynomials of
      Proposition~\ref{prop:strata} and the degree bounds $B_k$ of
      Proposition~\ref{prop:bound}.
\item \textbf{Emit.} For each $k$, output $A_k$: the tower of $A$ with the
      top element replaced by a generic linear ODE of order $k$ with
      polynomial coefficients of $v$-degree $\le B_k$, its coefficient list
      forming the new constants $c'$. Record the transfer maps
      $c \mapsto c'(c)$ (coefficients of $\widetilde{G}$ on each stratum).
\item \textbf{Run.} Apply the core algorithm --- projection plus the
      bad-locus recursion of the paper's general algorithm --- to $A$ and to
      each $A_k$; return the union of the membership loci (and, if the
      original-space geometry is wanted, the strata of step 2).
\end{enumerate}

\begin{theorem}[covering]
\label{thm:covering}
Let $A$ be as in \S\ref{sec:setting} and $P$ linear of the given shape. For
every $c^*$ with $h(c^*) \neq 0$ (and satisfying the base-splitting
genericity of Remark~\ref{rem:splitting}) and every nonzero
$\Psi_0 \in \Sol(A(c^*)) \cap \Sol(P)$: either $c^* \in \Var(A)$, or,
with $k = \ord G(c^*) \in \{1, \dots, n-1\}$,
\[
\Psi_0 \in \Sol\big(A_k(c'(c^*))\big)
\qquad\text{and}\qquad
c'(c^*) \in \Var(A_k).
\]
Hence $\{A, A_1, \dots, A_{n-1}\}$ $\forall$-covers $A$ away from the bad
locus $\{h = 0\}$; parameter values on the bad locus are handled, exactly as
in the paper, by the general algorithm's recursion, with the present
construction applied to each ansatz that recursion produces (both recursions
strictly decrease their measures --- order here, the bad-locus dimension
there --- so the composition terminates).
\end{theorem}

\begin{proof}
Immediate from Theorems~\ref{thm:structure} and~\ref{thm:membership},
Proposition~\ref{prop:bound} guaranteeing that $\widetilde{G}$'s coefficients
fit within $A_k$'s degree budget.
\end{proof}

\section{The two-line example, closed}

The previous note's Example~1 exhibited the gap; the construction now closes
it mechanically. Take
\[
A(c):\ \ \Psi'' - c\,\Psi = 0
\qquad\qquad
P:\ \ \Psi' - \Psi = 0 ,
\]
univariate base, so $L = \partial^2 - c$ and the remainder is the single
operator $R = \partial - 1$ ($n = 2$, $n' = 1$, $D = 0$).

\emph{Step 1--2.} One right division,
\[
\partial^2 - c \;=\; (\partial + 1)(\partial - 1) \;+\; (1 - c),
\]
shows the $0$-th subresultant (the differential resultant) is $1 - c$ up to a
unit. Strata: for $c \neq 1$, $G = 1$ --- no common solutions; on
$\{c = 1\}$, $G = \partial - 1$, order $k = 1$. The partial stratum of the
previous note is recovered as the vanishing locus of one polynomial in $c$,
with no decomposition performed.

\emph{Step 3.} $B_1 = (n + n' - 2k)D = 0$: emit
$A_1(c'):\ \Psi' - c'\Psi = 0$, constant coefficient. Transfer map:
$c'(c) = 1$ on the stratum.

\emph{Step 4.} The core algorithm on $A_1$ reduces $P$ to
$r_1 = (c' - 1)\Psi$, so $\Var(A_1) = \{c' = 1\}$ --- the family
$\Psi = K e^{x}$, containing the witness that was invisible to $\Var(A) =
\varnothing$. The previous note's \S4 family, introduced there by hand, is
here \emph{derived}, together with the stratum equation $c = 1$ and the point
$c' = 1$ it maps to.

\section{Beyond a single linear element}
\label{sec:beyond}

\paragraph{Towers (nested ODE elements).}
For ans\"atze like the paper's Figure ``ansatz 12'', a witness may instead
constrain the integration constants of a \emph{lower} level (a special
$\Theta$). The construction recurses: the subresultant coefficients from
step~2, which now live in the lower tower's differential field rather than in
$K(v)$, become the constraint one level down; treat that level as top and
repeat. GCRD theory is valid over any differential field, so
Theorems~\ref{thm:structure} and~\ref{thm:membership} survive level by
level; what must be re-proved per level is the \emph{uniform degree bound}
(the analogue of Proposition~\ref{prop:bound} over the lower tower's fraction
field). That lemma is of the same character as the bounded-prolongation
note's $N_0$ bound and is left open here. Orders strictly decrease and the
tower is finite, so the recursion terminates in a finite collection
regardless.

\paragraph{Nonlinear first-order elements.}
The vector-space machinery is linear to its bones, but the first-order
nonlinear case has its own closure. Full reduction against a first-order
element leaves a remainder containing \emph{no} $\Psi'$ at all, so a
partial-stratum witness satisfies the \emph{algebraic} relation
$r(v, \Psi) = 0$; the refined system replaces the ODE element by an
irreducible factor $q$ of $r$ as an \emph{algebraic extension element} (a
``green box''), of degree $\le \deg_\Psi r$ --- a bound available from the
reduction. Compatibility of $q$ with the derivation becomes polynomial
conditions on the constants, playing the role of the transfer map. So the
closure of a Riccati-type element consists of algebraic-extension ans\"atze
of bounded degree, again entirely within the paper's formalism.

\paragraph{Higher-order nonlinear elements.}
Open. Witnesses may sit on singular components (separant loci), where
neither the GCRD argument nor the algebraic-factor argument applies; Ritt's
theory of singular solutions is the natural starting point, and no
bounded-shape theorem is claimed.

\section{Reading}

The membership algorithm's blindness to the partial-solution strata is not a
price of its cheapness --- for linear ODE ans\"atze it is a bookkeeping gap,
closed by enlarging the ansatz \emph{collection} rather than the
\emph{algorithm}. Everything added is linear algebra over $K(v)$ with entries
polynomial in the constants: no differential elimination is performed beyond
the single Ritt reduction the core algorithm already does. The existential
formulation's extra content --- which subvarieties of constant space carry
partial solutions, and what those solutions are --- reappears as the
subresultant stratification (Proposition~\ref{prop:strata}) and the
lower-order membership families (Theorem~\ref{thm:covering}).

The case $k = 1$ connects to classical territory: a one-dimensional common
kernel is an exponential solution, i.e.\ a first-order right factor, the
object computed by the Beke and van~Hoeij factorization algorithms
\cite{beke,vanhoeij}. Their local (Newton-polygon) degree bounds are sharper
than Proposition~\ref{prop:bound}'s determinant bound, but are pointwise in
$c$; the subresultant bound is cruder and uniform, which is what a
parametric family construction requires. Where the sharper bounds matter in
practice, they can be substituted stratum by stratum.

\begin{thebibliography}{9}

\bibitem{ore} O.~Ore,
\emph{Theory of non-commutative polynomials},
Ann.\ of Math.\ \textbf{34} (1933), 480--508.

\bibitem{chardin} M.~Chardin,
\emph{Differential resultants and subresultants},
Proc.\ Fundamentals of Computation Theory (FCT'91),
Lecture Notes in Comput.\ Sci.\ \textbf{529}, Springer, 1991, 180--189.

\bibitem{li} Z.~Li,
\emph{A subresultant theory for Ore polynomials with applications},
Proc.\ ISSAC~'98, ACM, 1998, 132--139.

\bibitem{vdps} M.~van der Put and M.~F.~Singer,
\emph{Galois Theory of Linear Differential Equations},
Grundlehren der mathematischen Wissenschaften \textbf{328},
Springer, 2003.

\bibitem{beke} E.~Beke,
\emph{Die Irreducibilit\"at der homogenen linearen Differentialgleichungen},
Math.\ Ann.\ \textbf{45} (1894), 278--294.

\bibitem{vanhoeij} M.~van Hoeij,
\emph{Factorization of differential operators with rational functions
coefficients},
J.\ Symbolic Comput.\ \textbf{24} (1997), 537--561.

\end{thebibliography}

\end{document}
